\SourceLeaf{647}
\section*{CONTENTS.}
\addcontentsline{toc}{section}{Contents}
\begin{center}\bfseries GAUSS’S WORKS, VOLUME V.\par MATHEMATICAL PHYSICS.\end{center}
\par\nopagebreak[4]\smallskip\noindent\makebox[\linewidth][c]{\rule{.075\linewidth}{.35pt}}\par
\begingroup\fontsize{9}{11}\selectfont\raggedright
\par\Needspace{5\baselineskip}\medskip\noindent\textit{Papers}\par\smallskip
\noindent\hangindent=1em Theory of the Attraction of Homogeneous Ellipsoidal Bodies \unskip\hspace{1em plus 1fill}\mbox{1813 March . .\quad Page 1}\par\smallskip
\noindent\hangindent=1em On a New General Fundamental Law of Mechanics \unskip\hspace{1em plus 1fill}\mbox{1829\quad — 23}\par\smallskip
\noindent\hangindent=1em General Principles of the Theory of the Figure of Fluids in Equilibrium . . .\unskip\hspace{1em plus 1fill}\mbox{1829 Sept. . .\quad — 29}\par\smallskip
\noindent\hangindent=1em The Intensity of the Earth’s Magnetic Force Reduced to Absolute Measure . . .\unskip\hspace{1em plus 1fill}\mbox{1832 Dec. . .\quad — 79}\par\smallskip
\noindent\hangindent=1em General Theory of Terrestrial Magnetism . . . .\unskip\hspace{1em plus 1fill}\mbox{1838\quad 119}\par\smallskip
\noindent\hangindent=1em General Theorems on Attractive and Repulsive Forces That Act Inversely as the Square of the Distance \unskip\hspace{1em plus 1fill}\mbox{1839\quad — 195}\par\smallskip
\noindent\hangindent=1em Dioptrical Investigations . . . . . . . . . . . . . .\unskip\hspace{1em plus 1fill}\mbox{1840 Dec. . .\quad — 243}\par\smallskip
\par\Needspace{5\baselineskip}\medskip\noindent\textit{Notices of the Author’s Own Papers.}\par\smallskip
\noindent\hangindent=1em Theory of the Attraction of Homogeneous Ellipsoidal Bodies .\unskip\hspace{1em plus 1fill}\mbox{1813 April . .\quad — 279}\par\smallskip
\noindent\hangindent=1em General Principles of the Theory of the Figure of Fluids in Equilibrium . . . .\unskip\hspace{1em plus 1fill}\mbox{1829 Oct. . .\quad — 287}\par\smallskip
\noindent\hangindent=1em The Intensity of the Earth’s Magnetic Force Reduced to Absolute Measure . . .\unskip\hspace{1em plus 1fill}\mbox{1832 Dec. . .\quad — 293}\par\smallskip
\noindent\hangindent=1em General Theorems Relating to the Inverse Ratio, etc. \unskip\hspace{1em plus 1fill}\mbox{1840 March .\quad — 305}\par\smallskip
\noindent\hangindent=1em Dioptrical Investigations . . . . . . . . . . . . . . .\unskip\hspace{1em plus 1fill}\mbox{1841 Jan. . .\quad — 309}\par\smallskip
\par\Needspace{5\baselineskip}\medskip\noindent\textit{Miscellaneous Papers on Magnetism.}\par\smallskip
\noindent\hangindent=1em Terrestrial Magnetism and Magnetometers . . . . . . . . . .\unskip\hspace{1em plus 1fill}\mbox{1836\quad — 315}\par\smallskip
\noindent\hangindent=1em Introduction to the Journal Results, etc. . . . . . . . .\unskip\hspace{1em plus 1fill}\mbox{1836\quad — 345}\par\smallskip
\noindent\hangindent=1em A New Aid for Magnetic Observations . . . . . .\unskip\hspace{1em plus 1fill}\mbox{1837 Oct\quad — 352}\par\smallskip
\noindent\hangindent=1em On a New Instrument Intended Primarily for Direct Observation of Changes in the Intensity of the Horizontal Component of Terrestrial Magnetism . . . . . . . . . . . . . . . . .\unskip\hspace{1em plus 1fill}\mbox{1837\quad — 357}\par\smallskip
\noindent\hangindent=1em Instructions for Determining the Period of Oscillation of a Magnetic Needle . .\unskip\hspace{1em plus 1fill}\mbox{1837\quad — 374}\par\smallskip
\noindent\hangindent=1em On a Way to Make Deflections Easier to Observe . .\unskip\hspace{1em plus 1fill}\mbox{1839\quad — 395}\par\smallskip
\noindent\hangindent=1em On Determining the Constants of the Bifilar Magnetometer . . . . .\unskip\hspace{1em plus 1fill}\mbox{1840\quad — 404}\par\smallskip
\noindent\hangindent=1em Instructions for Determining the Magnetic Effect of a Bar Magnet at a Distance . . . . . . . . . . . . . .\unskip\hspace{1em plus 1fill}\mbox{1840\quad — 427}\par\smallskip
\SourceLeaf{648}
\noindent\hangindent=1em On the Use of the Magnetometer to Determine Absolute Declination . . . . . . . . . . . . . . . . . . . .\unskip\hspace{1em plus 1fill}\mbox{1841\quad Page 436}\par\smallskip
\noindent\hangindent=1em Observations of Magnetic Inclination in Göttingen . . . . . .\unskip\hspace{1em plus 1fill}\mbox{1841\quad — 444}\par\smallskip
\par\Needspace{5\baselineskip}\medskip\noindent\textit{Papers on Various Topics in Mathematical Physics.}\par\smallskip
\noindent\hangindent=1em Fundamental Equations for the Motion of Heavy Bodies on the Rotating Earth . . . . . . . . . . . . . . . . . . . .\unskip\hspace{1em plus 1fill}\mbox{1803\quad — 495}\par\smallskip
\noindent\hangindent=1em On Achromatic Doublet Objectives, Especially with a View to More Completely Eliminating Chromatic Dispersion . . . . . . . .\unskip\hspace{1em plus 1fill}\mbox{1817 Dec. . .\quad — 504}\par\smallskip
\noindent\hangindent=1em Letter to \textsc{Brandes} on the Same Subject . . . . . . . . . .\unskip\hspace{1em plus 1fill}\mbox{1831\quad — 509}\par\smallskip
\noindent\hangindent=1em Adjusting the Position of the Knife Edges of a Balance . . . . . . .\unskip\hspace{1em plus 1fill}\mbox{1837 March . .\quad — 511}\par\smallskip
\par\Needspace{5\baselineskip}\medskip\noindent\textit{Physical Observations.}\par\smallskip
\noindent\hangindent=1em Northern Lights on 7 January 1831 . . . . . . . . . . . . . . .\unskip\hspace{1em plus 1fill}\mbox{1831 Feb. . .\quad — 517}\par\smallskip
\noindent\hangindent=1em Magnetic Observatory in Göttingen . . . . . . . . . . .\unskip\hspace{1em plus 1fill}\mbox{1834 Aug. . .\quad — 519}\par\smallskip
\noindent\hangindent=1em Observations of Magnetic Variation in Göttingen and Leipzig . .\unskip\hspace{1em plus 1fill}\mbox{1834 Oct. . .\quad — 523}\par\smallskip
\noindent\hangindent=1em Magnetic Observatory in Göttingen . . . . . . . . . . .\unskip\hspace{1em plus 1fill}\mbox{1835 March . .\quad — 528}\par\smallskip
\noindent\hangindent=1em Observations of Magnetic Variation in Copenhagen and Milan .\unskip\hspace{1em plus 1fill}\mbox{1835 March . .\quad — 537}\par\smallskip
\noindent\hangindent=1em Magnetic Observatory in Göttingen . . . . . . . . . . .\unskip\hspace{1em plus 1fill}\mbox{1836 June . .\quad — 540}\par\smallskip
\noindent\hangindent=1em Procedure to Be Followed at the Scheduled Observations . . . . .\unskip\hspace{1em plus 1fill}\mbox{1836\quad — 541}\par\smallskip
\noindent\hangindent=1em Extract from Three Years of Daily Observations of Magnetic Declination in Göttingen . . . . . . . . . . . . . . . . .\unskip\hspace{1em plus 1fill}\mbox{1836\quad — 556}\par\smallskip
\noindent\hangindent=1em Explanations of the Plots and Numerical Readings for the Scheduled Observations .\unskip\hspace{1em plus 1fill}\mbox{1836\quad — 568}\par\smallskip
\noindent\hangindent=1em Explanations of the Plots and Numerical Readings for the Scheduled Observations .\unskip\hspace{1em plus 1fill}\mbox{1837\quad — 576}\par\smallskip
\noindent\hangindent=1em The Earth’s Magnetic South Pole . . . . . . . . . . . . .\unskip\hspace{1em plus 1fill}\mbox{1841\quad — 580}\par\smallskip
\par\Needspace{5\baselineskip}\medskip\noindent\textit{Reviews of Works by Other Authors.}\par\smallskip
\noindent\hangindent=1em \textsc{Benzenberg.} On \textsc{Dalton}'s Theory . . . . . . . . . .\unskip\hspace{1em plus 1fill}\mbox{1830 Dec. . .\quad — 583}\par\smallskip
\noindent\hangindent=1em \textsc{Fischer.} Artificial Magnets . . . . . . . . . . . . . . .\unskip\hspace{1em plus 1fill}\mbox{1832 Sept. .\quad — 591}\par\smallskip
\noindent\hangindent=1em Results of the Observations of the Magnetic Association . . . . .\unskip\hspace{1em plus 1fill}\mbox{1837 June . .\quad — 595}\par\smallskip
\noindent\hangindent=1em \textsc{Gerling.} Organization of the Mathematical and Physical Institute . . .\unskip\hspace{1em plus 1fill}\mbox{1842 April .\quad — 596}\par\smallskip
\par\Needspace{17\baselineskip}\medskip\noindent\textit{Posthumous Papers.}\par\smallskip
\noindent\hangindent=1em On Electrodynamics . . . . . . . . . . . . . . . . . . . . . . . .\unskip\hspace{1em plus 1fill}\mbox{\quad — 601}\par\smallskip
\noindent\hangindent=1em On Spherical Harmonics . . . . . . . . . . . . . . . . . . . . . . . .\unskip\hspace{1em plus 1fill}\mbox{\quad — 630}\par\smallskip
\noindent\hangindent=1em On Using the Comparator . . . . . . . . . . . . . . . . . . . .\unskip\hspace{1em plus 1fill}\mbox{\quad — 632}\par\smallskip
\noindent\hangindent=1em General Formulas for the Effect of a Luminous Point $P$ on a Point $p$ . .\unskip\hspace{1em plus 1fill}\mbox{\quad — 635}\par\smallskip
\noindent\hangindent=1em Remarks . . . . . . . . . . . . . . . . . . . . . . . . . . . .\unskip\hspace{1em plus 1fill}\mbox{\quad — 637}\par\smallskip
\par\medskip\noindent Lithographic plate for the theory of terrestrial magnetism, page 176.
\endgroup

\par\nopagebreak[4]\medskip\noindent\makebox[\linewidth][c]{\rule{.15\linewidth}{.35pt}}\par\nopagebreak[4]
\bigskip
\begin{center}
GÖTTINGEN,\par
\small PRINTED AT DIETERICH’S UNIVERSITY PRESS\par
W. FR. KAESTNER.
\end{center}
